\section{Introduction}
\label{sec:introduction}

This is the fourth and final paper in the series.
Paper~I established the substrate model.
Papers~II and~III developed the Lorentz, gravity, gauge, and color bridges.
The present paper develops two bridges that have been deliberately held back: matter and mass.

\paragraph{Matter bridge.}
The matter bridge poses particle-like states as self-adjoint eigenproblems on a symmetric
background, with angular structure organized by vector spherical harmonics (\Cref{sec:matter}).
The key structural inputs are the geometric quartic $W_4\propto k_s\alpha/a$ from Paper~I and
the $U(3)$ triplet structure from Paper~III.
Two mechanisms together establish stability:
(i) anti-collapse is provided by the lattice spacing $a$ --- the Derrick $\lambda\to 0$
rescaling halts at one cell, because configurations concentrated below one cell are not local
energy minima of the spring network;
(ii) anti-expansion is provided by the geometric quartic $\propto k_s\alpha/a$, which under
Derrick rescaling scales as $\lambda^{-1}$ and overpowers the $\lambda^{+1}$ gradient term for
$\lambda > \lambda^*$.
The equilibrium soliton radius from this balance is $R_h/a \approx \kappa\,(A/a)
\sqrt{\alpha/(1-\alpha)}$, setting a soliton sweet spot at $\alpha\approx 0.5$--$0.8$.

\paragraph{Mass bridge.}
The mass bridge derives emergent rest mass from self-confinement of a closed-loop wave mode ---
not from the Higgs mechanism (\Cref{sec:mass}).
A wave trapped in a circularly closed toroidal excitation cannot disperse: it is pinned by topology
and geometric self-guidance.
The total substrate energy of this mode at equilibrium, $E(R^*,A)$, is identified with the rest
mass $mc^2$.
The complex structure $\Psi\approx\xi+(i/\omega_0)\dot\xi$ that mediates this picture arises
entirely from the timelike worldtube direction --- the imaginary unit $i$ here is a
quarter-period time shift, not an independent complex field.
A separate, quantitative open problem is connecting the equilibrium radius $R^*$ to the Compton
wavelength $\hbar/mc$; this requires deriving $\hbar$ and $c$ from substrate parameters
consistently.

\paragraph{Why matter and mass are paired.}
Both bridges are gated on the same object: a converged, stable, localized soliton eigenstate.
Without such an eigenstate, the mass identification $E(R^*,A)=mc^2$ remains schematic.
The matter bridge's decisive numerical task (does the substrate support a stable baryon-like
$\pi_3$ texture?) is therefore simultaneously the gate that unlocks the mass bridge's
normalization argument.
Pairing them in one paper makes this dependency explicit.

\paragraph{What this paper does not claim.}
Which topological objects the substrate dynamics actually stabilizes is an open numerical question.
This paper does \emph{not} claim the identification of any specific topological object with a
specific physical particle; it develops the eigenproblem and stability criteria that such an
identification would have to satisfy, and poses the existence of a stable baryon-like texture as
its decisive numerical target.

\paragraph{Paper organization.}
\Cref{sec:interfaces} lists imports from Papers~I and~III.
\Cref{sec:matter} develops the matter bridge: self-adjoint eigenproblem, VSH angular structure,
Derrick stability, soliton labels, and periodic-BC requirement.
\Cref{sec:mass} develops the mass bridge: self-confinement picture, toroidal mode energy, Derrick
balance for torus geometry, time-link binding, and open normalization problems.